\clearpage
\fancyhead[L]{\small BESTIMMUNG DER GENAUIGKEIT DER BEOBACHTUNGEN}
\pdfbookmark[0]{Bestimmung der Genauigkeit der Beobachtungen}{genauigkeit}
\SourceLeaf{114}
\begin{center}
\rule{.8\linewidth}{.3pt}\par\smallskip
{\small Zeitschrift für Astronomie, herausgegeben von B. von \textsc{Lindenau} und J. G. F. \textsc{Bohnenberger}.\par
Band I. Heft für März und April 1816.\par}
\smallskip\rule{.8\linewidth}{.3pt}\par\bigskip
{\large\itshape Bestimmung der Genauigkeit der Beobachtungen.\par}
\end{center}
\renewcommand{\Article}[1]{\par\medskip\Needspace{4\baselineskip}\pdfbookmark[1]{Art. #1}{genauigkeit-art-#1}\centerline{\textbf{#1.}}\nobreak\smallskip}
\Article{1}
Bei der Begründung der sogenannten Methode der kleinsten Quadrate wird angenommen, dass sich die Wahrscheinlichkeit eines Beobachtungsfehlers \(\Delta\) durch die Formel
\[
\frac{h}{\surd\pi}.e^{-hh\Delta\Delta}
\]
ausdrückt, wobei \(\pi\) den halben Kreisumfang, \(e\) die Basis der hyperbolischen Logarithmen und \(h\) eine Konstante bezeichnet, die man nach Art. 178 der \textit{Theoria Motus Corporum Coelestium} als Maß für die Genauigkeit der Beobachtungen ansehen kann. Bei der Anwendung der Methode der kleinsten Quadrate zur Bestimmung der wahrscheinlichsten Werte der Größen, von denen die Beobachtungen abhängen, braucht man den Wert von \(h\) überhaupt nicht zu kennen; auch das \textit{Verhältnis} der Genauigkeit der Ergebnisse zu der Genauigkeit der Beobachtungen ist von \(h\) unabhängig. Dennoch ist die Kenntnis dieser Größe selbst stets interessant und lehrreich. Ich will daher zeigen, wie man durch die Beobachtungen selbst zu einer solchen Kenntnis gelangen kann.

\Article{2}
Zunächst stelle ich einige den Gegenstand erläuternde Bemerkungen voran. Der Kürze halber bezeichne ich den Wert des Integrals
\SourceLeaf{115}
\[
\int\frac{2e^{-tt}\,\mathrm{d}t}{\surd\pi}
\]
von \(t=0\) an gerechnet mit \(\theta t\). Einige einzelne Werte werden einen Eindruck vom Verlauf dieser Funktion geben. Man hat
\[
\begin{array}{r@{\;=\;}l@{\;=\;}l}
0{,}5000000&\theta\,0{,}4769363&\theta\rho\\
0{,}6000000&\theta\,0{,}5951161&\theta\,1{,}247790\rho\\
0{,}7000000&\theta\,0{,}7328691&\theta\,1{,}536618\rho\\
0{,}8000000&\theta\,0{,}9061939&\theta\,1{,}900032\rho\\
0{,}8427008&\theta\,1&\theta\,2{,}096716\rho\\
0{,}9000000&\theta\,1{,}1630872&\theta\,2{,}438664\rho\\
0{,}9900000&\theta\,1{,}8213864&\theta\,3{,}818930\rho\\
0{,}9990000&\theta\,2{,}3276754&\theta\,4{,}880475\rho\\
0{,}9999000&\theta\,2{,}7510654&\theta\,5{,}768204\rho\\
1&\multicolumn{2}{l}{\theta\infty}
\end{array}
\]

Die Wahrscheinlichkeit, dass der Fehler einer Beobachtung zwischen den Grenzen \(-\Delta\) und \(+\Delta\) liegt oder, ohne Rücksicht auf sein Vorzeichen, nicht größer als \(\Delta\) ist, beträgt
\[
=\int\frac{he^{-hhxx}\,\mathrm{d}x}{\surd\pi}
\]
wenn man das Integral von \(x=-\Delta\) bis \(x=+\Delta\) erstreckt; sie ist also doppelt so groß wie dasselbe Integral, wenn man es von \(x=0\) bis \(x=\Delta\) nimmt, mithin
\[
=\theta h\Delta
\]

Die Wahrscheinlichkeit, dass der Fehler nicht kleiner als \(\frac\rho h\) ist, beträgt folglich \(=\frac12\), also ebenso viel wie die Wahrscheinlichkeit des Gegenteils. Wir wollen diese Größe \(\frac\rho h\) den \textit{wahrscheinlichen Fehler} nennen und mit \(r\) bezeichnen. Dagegen beträgt die Wahrscheinlichkeit, dass der Fehler \(2{,}438664r\) übersteigt, nur \(\frac1{10}\); die Wahrscheinlichkeit, dass der Fehler \(3{,}818930r\) übersteigt, nur \(\frac1{100}\) usw.

\Article{3}
Nehmen wir nun an, dass bei \(m\) tatsächlich angestellten Beobachtungen die Fehler \(\alpha,\,\beta,\,\gamma,\,\delta\) usw. begangen worden sind, und untersuchen wir, was sich daraus hinsichtlich der Werte von \(h\) und \(r\) schließen lässt. Macht man zwei \SourceLeaf{116}%
Annahmen, indem man den wahren Wert von \(h\) einmal \(=H\) und einmal \(=H'\) setzt, so verhalten sich die Wahrscheinlichkeiten, mit denen man unter diesen Voraussetzungen die Fehler \(\alpha,\,\beta,\,\gamma,\,\delta\) usw. erwarten würde, zueinander wie

\[
\begin{gathered}
He^{-HH\alpha\alpha}.He^{-HH\beta\beta}.He^{-HH\gamma\gamma}.\text{ u.\,s.\,w.}\\
\text{zu}\quad H'e^{-H'H'\alpha\alpha}.H'e^{-H'H'\beta\beta}.H'e^{-H'H'\gamma\gamma}.\text{ u.\,s.\,w.}
\end{gathered}
\]
das heißt wie
\[
H^m e^{-HH(\alpha\alpha+\beta\beta+\gamma\gamma+\text{ u.\,s.\,w.})}\quad\text{zu}\quad
H'{}^m e^{-H'H'(\alpha\alpha+\beta\beta+\gamma\gamma+\text{ u.\,s.\,w.})}
\]
Im selben Verhältnis stehen folglich \textit{nach} dem Eintritt jener Fehler auch die Wahrscheinlichkeiten, dass \(H\) oder \(H'\) der wahre Wert von \(h\) war (\textit{T. M. C. C.}, Art. 176); oder die Wahrscheinlichkeit jedes Wertes von \(h\) ist der Größe
\[
h^m e^{-hh(\alpha\alpha+\beta\beta+\gamma\gamma+\text{ u.\,s.\,w.})}
\]
proportional. Der \textit{wahrscheinlichste} Wert von \(h\) ist folglich derjenige, für den diese Größe ein Maximum wird; nach den bekannten Regeln findet man ihn zu
\[
=\surd\frac{m}{2(\alpha\alpha+\beta\beta+\gamma\gamma+\text{ u.\,s.\,w.})}
\]
Der wahrscheinlichste Wert von \(r\) ist folglich
\[
\begin{aligned}
&=\rho\surd\frac{2(\alpha\alpha+\beta\beta+\gamma\gamma+\text{ u.\,s.\,w.})}{m}\\[4pt]
&=0{,}6744897\surd\frac{\alpha\alpha+\beta\beta+\gamma\gamma+\text{ u.\,s.\,w.}}{m}
\end{aligned}
\]
Dieses Ergebnis gilt allgemein, mag \(m\) groß oder klein sein.

\Article{4}
Es ist leicht einzusehen, dass man umso weniger berechtigt ist, von dieser Bestimmung von \(h\) und \(r\) eine große Genauigkeit zu erwarten, je kleiner \(m\) ist. Entwickeln wir daher den Genauigkeitsgrad, der dieser Bestimmung beizulegen ist, für den Fall, dass \(m\) eine große Zahl ist. Der Kürze halber bezeichnen wir den gefundenen wahrscheinlichen Wert von \(h\)
\[
\surd\frac{m}{2(\alpha\alpha+\beta\beta+\gamma\gamma+\text{ u.\,s.\,w.})}
\]
mit \(H\) und bemerken, dass sich die Wahrscheinlichkeit, \(H\) sei der
\SourceLeaf{117}%
wahre Wert von \(h\), zu der Wahrscheinlichkeit, der wahre Wert sei \(=H+\lambda\), verhält wie
\[
H^m e^{-\frac m2}:(H+\lambda)^m e^{-\frac{m(H+\lambda)^2}{2HH}}
\]
oder wie
\[
1:e^{-\frac{\lambda\lambda m}{HH}(1-\frac13.\frac\lambda H+\frac{1\,\lambda\lambda}{4HH}-\frac15.\frac{\lambda^3}{H^3}\, .\,.\,.\,.)}
\]
Das zweite Glied ist gegenüber dem ersten nur dann noch merklich, wenn \(\frac\lambda H\) ein kleiner Bruch ist. Daher dürfen wir uns erlauben, anstelle des angegebenen Verhältnisses dieses zu verwenden:
\[
1:e^{-\frac{\lambda\lambda m}{HH}}
\]
Dies bedeutet eigentlich: Die Wahrscheinlichkeit, dass der wahre Wert von \(h\) zwischen \(H+\lambda\) und \(H+\lambda+\mathrm{d}\lambda\) liegt, ist sehr nahe
\[
=Ke^{-\frac{\lambda\lambda m}{HH}}\mathrm{d}\lambda
\]
wobei \(K\) eine Konstante ist, die so bestimmt werden muss, dass das Integral
\[
\int Ke^{-\frac{\lambda\lambda m}{HH}}\mathrm{d}\lambda
\]
zwischen den zulässigen Grenzen von \(\lambda\) genommen gleich \(=1\) wird. Hier darf man statt solcher Grenzen, da wegen der Größe von \(m\)
\[
e^{-\frac{\lambda\lambda m}{HH}}
\]
offenbar unmerklich wird, sobald \(\frac\lambda H\) kein kleiner Bruch mehr ist, die Grenzen \(-\infty\) und \(+\infty\) nehmen; dadurch ergibt sich
\[
K=\frac1H\surd\frac m\pi
\]
Folglich ist die Wahrscheinlichkeit, dass der wahre Wert von \(h\) zwischen \(H-\lambda\) und \(H+\lambda\) liegt,
\[
=\theta\bigl(\tfrac\lambda H\surd m\bigr)
\]
also beträgt diese Wahrscheinlichkeit \(=\frac12\), wenn
\[
\frac\lambda H\surd m=\rho\quad\text{ist.}
\]
\SourceLeaf{118}%
Die Chancen stehen also eins zu eins dafür, dass der wahre Wert von \(h\)
\[
\text{zwischen}\quad H\bigl(1-\tfrac\rho{\surd m}\bigr)\quad\text{und}\quad H\bigl(1+\tfrac\rho{\surd m}\bigr)
\]
liegt oder dass der wahre Wert von \(r\)
\[
\text{zwischen}\quad\frac{R}{1-\frac\rho{\surd m}}\quad\text{und}\quad\frac{R}{1+\frac\rho{\surd m}}
\]
fällt, wobei \(R\) den im vorigen Artikel gefundenen wahrscheinlichsten Wert von \(r\) bezeichnet. Diese Grenzen kann man die \textit{wahrscheinlichen Grenzen der wahren Werte von \(h\) und \(r\)} nennen; offensichtlich können wir für die wahrscheinlichen Grenzen des wahren Wertes von \(r\) hier auch
\[
R\bigl(1-\tfrac\rho{\surd m}\bigr)\quad\text{und}\quad R\bigl(1+\tfrac\rho{\surd m}\bigr)
\]
setzen.

\Article{5}
In der vorangehenden Untersuchung sind wir davon ausgegangen, dass \(\alpha,\,\beta,\,\gamma,\,\delta\) usw. bestimmte, gegebene Größen seien, und haben die Wahrscheinlichkeit gesucht, dass der wahre Wert von \(h\) oder \(r\) zwischen gewissen Grenzen liegt. Man kann die Frage auch von einer anderen Seite betrachten: Unter der Voraussetzung, dass die Beobachtungsfehler einem bestimmten Wahrscheinlichkeitsgesetz folgen, lässt sich die Wahrscheinlichkeit bestimmen, mit der die Summe der Quadrate von \(m\) Beobachtungsfehlern zwischen gewisse Grenzen fällt. Für großes \(m\) hat \textsc{Laplace} diese Aufgabe bereits gelöst, ebenso wie diejenige, die Wahrscheinlichkeit dafür zu finden, dass die Summe von \(m\) Beobachtungsfehlern selbst zwischen gewisse Grenzen fällt. Diese Untersuchung lässt sich leicht weiter verallgemeinern; hier begnüge ich mich damit, das Ergebnis anzugeben.

Es bezeichne \(\varphi x\) die Wahrscheinlichkeit des Beobachtungsfehlers \(x\), sodass \(\int\varphi x.\mathrm{d}x=1\) gilt, wenn das Integral von \(x=-\infty\) bis \(x=+\infty\) erstreckt wird. Zwischen denselben Grenzen wollen wir allgemein den Wert des Integrals
\[
\int\varphi x.x^n\mathrm{d}x
\]
\SourceLeaf{119}
mit \(K^{\mathrm n}\) bezeichnen.\footnote{Oder vielmehr soll das Integral \(\int\varphi x.x^n\mathrm{d}x\) zwischen den Grenzen \(x=0\) und \(x=\infty\) mit
\[\tfrac12K^{\mathrm n}\]
bezeichnet werden. [Handschriftliche Bemerkung]} Ferner sei \(S^{\mathrm n}\) die Summe
\[
\alpha^n+\beta^n+\gamma^n+\delta^n+\text{ u.\,s.\,w.}
\]
wobei \(\alpha,\,\beta,\,\gamma,\,\delta\) usw. für \(m\) beliebige Beobachtungsfehler stehen; sämtliche Glieder dieser Summe sollen auch bei ungeradem \(n\) positiv genommen werden.

Dann ist \(mK^{\mathrm n}\) der wahrscheinlichste Wert von \(S^{\mathrm n}\), und die Wahrscheinlichkeit dafür, dass der wahre Wert von \(S^{\mathrm n}\) zwischen den Grenzen \(mK^{\mathrm n}-\lambda\) und \(mK^{\mathrm n}+\lambda\) liegt, beträgt
\[
=\theta\frac{\lambda}{\surd[2m(K^{2\mathrm n}-K^{\mathrm n}K^{\mathrm n})]}
\]
Folglich sind die wahrscheinlichen Grenzen von \(S^{\mathrm n}\)
\[
mK^{\mathrm n}-\rho\surd[2m(K^{2\mathrm n}-K^{\mathrm n}K^{\mathrm n})]\quad\text{und}\quad
mK^{\mathrm n}+\rho\surd[2m(K^{2\mathrm n}-K^{\mathrm n}K^{\mathrm n})]
\]
Dieses Ergebnis gilt allgemein für jedes Gesetz der Beobachtungsfehler. Wenden wir es auf den Fall an, dass
\[
\varphi x=\frac h{\surd\pi}e^{-hhxx}
\]
gesetzt wird, so erhalten wir
\[
K^{\mathrm n}=\frac{\Pi\,\frac12(n-1)}{h^n\surd\pi}
\]
wobei die Charakteristik \(\Pi\) im Sinne der \textit{Disquisitiones generales circa seriem infinitam} (Comm. nov. soc. Gotting. T. II.), M. 5, Art. 28 der angeführten Abhandlung, genommen ist. Also
\[
\begin{aligned}
K&=1,&K'&=\frac1{h\surd\pi},&K''&=\frac1{2hh},&K'''&=\frac1{h^3\surd\pi}\\[5pt]
K^{\mathrm{IV}}&=\frac{1.3}{4h^4},&K^{\mathrm{V}}&=\frac{1.2}{h^5\surd\pi},&K^{\mathrm{VI}}&=\frac{1.3.5}{8h^6},&K^{\mathrm{VII}}&=\frac{1.2.3}{h^7\surd\pi}\text{ u.\,s.\,w.}
\end{aligned}
\]
Der wahrscheinlichste Wert von \(S^{\mathrm n}\) ist demnach
\[
\frac{m\Pi\,\frac12(n-1)}{h^n\surd\pi}
\]
und die wahrscheinlichen Grenzen des wahren Wertes von \(S^{\mathrm n}\) sind
\SourceLeaf{120}
\[
\frac{m\Pi\,\frac12(n-1)}{h^n\surd\pi}\Bigl\{1-\rho\surd\Bigl(\frac2m.\Bigl(\frac{\Pi(n-\frac12).\surd\pi}{(\Pi\,\frac12(n-1))^2}-1\Bigr)\Bigr)\Bigr\}
\]
und
\[
\frac{m\Pi\,\frac12(n-1)}{h^n\surd\pi}\Bigl\{1+\rho\surd\Bigl(\frac2m.\Bigl(\frac{\Pi(n-\frac12).\surd\pi}{(\Pi\,\frac12(n-1))^2}-1\Bigr)\Bigr)\Bigr\}
\]
Setzt man also wie oben
\[
\frac\rho h=r
\]
sodass \(r\) den wahrscheinlichen Beobachtungsfehler bezeichnet, so ist der wahrscheinlichste Wert von
\[
\rho\BareRoot[n]{\frac{S^{\mathrm n}\surd\pi}{m\Pi\,\frac12(n-1)}}
\]
offenbar \(=r\); die wahrscheinlichen Grenzen des Wertes dieser Größe sind
\[
r\Bigl\{1-\frac\rho n\surd\Bigl(\frac2m.\Bigl(\frac{\Pi(n-\frac12).\surd\pi}{(\Pi\,\frac12(n-1))^2}-1\Bigr)\Bigr)\Bigr\}
\]
und
\[
r\Bigl\{1+\frac\rho n\surd\Bigl(\frac2m.\Bigl(\frac{\Pi(n-\frac12).\surd\pi}{(\Pi\,\frac12(n-1))^2}-1\Bigr)\Bigr)\Bigr\}
\]
Somit kann man ebenfalls mit gleichen Chancen wetten, dass \(r\) zwischen den Grenzen
\[
\begin{aligned}
&\rho\BareRoot[n]{\frac{S^{\mathrm n}\surd\pi}{m\Pi\,\frac12(n-1)}}\\[3pt]
&\quad\Bigl\{1-\frac\rho n\surd\Bigl(\frac2m.\Bigl(\frac{\Pi(n-\frac12).\surd\pi}{(\Pi\,\frac12(n-1))^2}-1\Bigr)\Bigr)\Bigr\}
\end{aligned}
\]
und
\[
\begin{aligned}
&\rho\BareRoot[n]{\frac{S^{\mathrm n}\surd\pi}{m\Pi\,\frac12(n-1)}}\\[3pt]
&\quad\Bigl\{1+\frac\rho n\surd\Bigl(\frac2m.\Bigl(\frac{\Pi(n-\frac12).\surd\pi}{(\Pi\,\frac12(n-1))^2}-1\Bigr)\Bigr)\Bigr\}
\end{aligned}
\]
liegt. Für \(n=2\) lauten diese Grenzen
\[
\rho\surd\frac{2S''}{m}\bigl\{1-\tfrac\rho{\surd m}\bigr\}\quad\text{und}\quad
\rho\surd\frac{2S''}{m}\bigl\{1+\tfrac\rho{\surd m}\bigr\}
\]
und stimmen vollständig mit den oben (Art. 4) gefundenen überein. Allgemein hat man für gerades \(n\) die Grenzen
\[
\begin{aligned}
&\rho\surd2.\BareRoot[n]{\frac{S^{\mathrm n}}{m.1.3.5.7\, .\,.\,.\,.\,(n-1)}}\\[4pt]
&\quad\Bigl\{1-\frac\rho n\surd\Bigl(\frac2m.\Bigl(\frac{(n+1).(n+3)\, .\,.\,.\,(2n-1)}{1.3.5\, .\,.\,.\,(n-1)}-1\Bigr)\Bigr)\Bigr\}
\end{aligned}
\]
und
\[
\begin{aligned}
&\rho\surd2.\BareRoot[n]{\frac{S^{\mathrm n}}{m.1.3.5.7\, .\,.\,.\,.\,(n-1)}}\\[4pt]
&\quad\Bigl\{1+\frac\rho n\surd\Bigl(\frac2m.\Bigl(\frac{(n+1).(n+3)\, .\,.\,.\,(2n-1)}{1.3.5\, .\,.\,.\,(n-1)}-1\Bigr)\Bigr)\Bigr\}
\end{aligned}
\]
und für ungerades \(n\) folgende:
\[
\begin{aligned}
&\rho\BareRoot[n]{\frac{S^{\mathrm n}\surd\pi}{m.1.2.3\, .\,.\,.\,\frac12(n-1)}}\\[4pt]
&\quad\Bigl\{1-\frac\rho n\surd\Bigl(\frac1m.\Bigl(\frac{1.3.5.7\, .\,.\,.\,(2n-1)\pi}{(2.4.6\, .\,.\,.\,(n-1))^2}-2\Bigr)\Bigr)\Bigr\}
\end{aligned}
\]
und
\[
\begin{aligned}
&\rho\BareRoot[n]{\frac{S^{\mathrm n}\surd\pi}{m.1.2.3\, .\,.\,.\,\frac12(n-1)}}\\[4pt]
&\quad\Bigl\{1+\frac\rho n\surd\Bigl(\frac1m.\Bigl(\frac{1.3.5.7\, .\,.\,.\,(2n-1)\pi}{(2.4.6\, .\,.\,.\,(n-1))^2}-2\Bigr)\Bigr)\Bigr\}
\end{aligned}
\]

\SourceLeaf{121}
\Article{6}
Ich füge noch die Zahlenwerte für die einfachsten Fälle hinzu:
\begin{center}Wahrscheinliche Grenzen von \(r\)\end{center}
\[
\begin{array}{r@{\quad}l}
\mathrm{I}.&0{,}8453473\quad\dfrac{S'}m.\bigl(1\mp\dfrac{0{,}5095841}{\surd m}\bigr)\\[10pt]
\mathrm{II}.&0{,}6744897\BareRoot[2]{\dfrac{S''}m}.\bigl(1\mp\dfrac{0{,}4769363}{\surd m}\bigr)\\[10pt]
\mathrm{III}.&0{,}5771897\BareRoot[3]{\dfrac{S'''}m}.\bigl(1\mp\dfrac{0{,}4971987}{\surd m}\bigr)\\[10pt]
\mathrm{IV}.&0{,}5125017\BareRoot[4]{\dfrac{S''''}m}.\bigl(1\mp\dfrac{0{,}5507186}{\surd m}\bigr)\\[10pt]
\mathrm{V}.&0{,}4655532\BareRoot[5]{\dfrac{S^{\mathrm V}}m}.\bigl(1\mp\dfrac{0{,}6355080}{\surd m}\bigr)\\[10pt]
\mathrm{VI}.&0{,}4294972\BareRoot[6]{\dfrac{S^{\mathrm{VI}}}m}.\bigl(1\mp\dfrac{0{,}7557764}{\surd m}\bigr)
\end{array}
\]
Man sieht auch hier, dass Verfahren II von allen am vorteilhaftesten ist. Hundert Beobachtungsfehler, nach dieser Formel behandelt, ergeben nämlich ein ebenso zuverlässiges Resultat wie
\begin{center}
114 nach I,\quad109 nach III,\quad133 nach IV,\quad178 nach V,\quad251 nach VI.
\end{center}
Die Formel I hat jedoch den Vorzug der mit Abstand bequemsten Rechnung. Da sie nicht viel ungenauer als II ist, kann man sie verwenden, wenn man die Summe der Fehlerquadrate weder ohnehin schon kennt noch kennen möchte.

\Article{7}
Noch bequemer, wenn auch beträchtlich ungenauer, ist folgendes Verfahren. Man ordne sämtliche \(m\) Beobachtungsfehler ihrem absoluten Betrag nach und nenne den mittleren Wert \(M\), wenn ihre Anzahl ungerade ist, oder bei gerader Anzahl das arithmetische Mittel der beiden mittleren Werte. Es lässt sich zeigen—doch kann dies hier nicht weiter ausgeführt werden—, dass bei einer großen Zahl von Beobachtungen \(r\) der wahrscheinlichste Wert von \(M\) ist und die wahrscheinlichen Grenzen von \(M\)
\[
r\bigl(1-e^{\rho\rho}\surd\tfrac\pi{8m}\bigr)\quad\text{und}\quad r\bigl(1+e^{\rho\rho}\surd\tfrac\pi{8m}\bigr)
\]
betragen; oder die wahrscheinlichen Grenzen des Wertes von \(r\) sind
\SourceLeaf{122}
\[
\begin{gathered}
M\bigl(1-e^{\rho\rho}\surd\tfrac\pi{8m}\bigr)\quad\text{und}\quad M\bigl(1+e^{\rho\rho}\surd\tfrac\pi{8m}\bigr),\\[4pt]
\text{oder in Zahlen}\quad M\bigl(1\mp\tfrac{0{,}7520974}{\surd m}\bigr)
\end{gathered}
\]
Dieses Verfahren ist also nur wenig genauer als die Anwendung der Formel VI; man müsste 249 Beobachtungsfehler heranziehen, um ebenso weit zu kommen wie mit 100 Beobachtungsfehlern nach Formel II.

\Article{8}
Die Anwendung einiger dieser Methoden auf die Fehler aus \textsc{Bode}s astronomischem Jahrbuch für 1818, S. 234, bei 48 Beobachtungen der Rektaszensionen des Polarsterns durch \textsc{Bessel} ergab
\[
S'=60''46,\quad S''=110''600,\quad S'''=250''341118
\]
Daraus ergaben sich die wahrscheinlichsten Werte von \(r\):
\begin{center}
\begin{tabular}{r@{\;}c@{\;}l@{\quad}l@{\;}l}
nach Formel I&.\,.\,.\,.&\(1''065,\)&wahrscheinliche Unsicherheit \(=\)&\(\pm0''078\)\\
II&.\,.\,.\,.&\(1{,}024\)&&\(\pm0{,}070\)\\
III&.\,.\,.\,.&\(1{,}001\)&&\(\pm0{,}072\)\\
nach Art. 7&.\,.\,.\,.\,.&\(1{,}045\)&&\(\pm0{,}113\)
\end{tabular}
\end{center}
eine Übereinstimmung, wie sie kaum zu erwarten war. \textsc{Bessel} selbst gibt \(1''067\) an und scheint daher nach Formel I gerechnet zu haben.
\begin{center}\rule{.24\linewidth}{.3pt}\end{center}
